\begin{abstract}
Tree speculative decoding for recurrent-hybrid language models must preserve candidate ancestry and the state of the selected continuation. We study GDN Tree-Scan, a Gated DeltaNet verifier integrated with vLLM and FlashAttention-2. Its accepted-prefix contract coordinates recurrent state, convolution history, attention key--value storage, and a pending correction token. Archived regression witnesses expose incomplete state publication and numerical disagreement caused by physical attention layout. The publication repair changes a token fixture from 15 failures in 15 checks to zero; the layout repair yields shared-spine byte agreement in its tested route. Greater acceptance does not guarantee speed: an archived batch-four-configured Qwen3.6-27B-FP8 campaign on GB10 commits 5.286 tokens per verification event versus a native five-step chain's 4.422, but achieves 32.85 versus 42.74 decode-step tokens/s. Fresh same-input measurements compare sequential scan/replay with local triangular-solve and finite-Neumann reconstructions on eight pilot prefixes and 31 of 32 confirmation prefixes passing frozen provenance checks. With fp32 products, errors against a float64 oracle are on the order of fp32 roundoff; tf32 increases error by roughly three orders of magnitude. Nevertheless, frozen bitwise, ordering, padding, and timing criteria fail. A retrospective audit also identifies omitted checks and inaccurate pilot rationale in the frozen record. Historical sampling and binary provenance limit the archived comparison; current model-level continuation and matched final-stack timing remain unqualified. The results provide scoped state-integrity witnesses and numerical/cost characterization, without a superiority or distribution-equivalence claim.
\end{abstract}
